\section{Capture-avoiding substitution}
\label{sec:fol-syntax-semantics-04}

\begin{definition}[Substitution in a term]
For terms \(r,t\) and a variable \(x\), the term \(r[t/x]\) is defined by
\[
 x[t/x]=t,\qquad y[t/x]=y\ (y\ne x),
\]
and
\[
 f(r_1,\ldots,r_n)[t/x]
 =f(r_1[t/x],\ldots,r_n[t/x]).
\]
Only occurrences of \(x\) are replaced.
\end{definition}

Fix the order \(x_0,x_1,\ldots\) on \(\Var\).  This order makes every choice
of a fresh variable below deterministic.

\begin{definition}[Capture-avoiding substitution in a formula]
For a formula \(\varphi\), term \(t\), and variable \(x\), the formula
\(\varphi[t/x]\) is defined by recursion on formula complexity.  In atomic
formulas, substitute in every term.  Substitution commutes with
\(\bottom,\neg,\land,\lor,\to\).  For
\(Q\in\set{\forall,\exists}\), define
\[
(Qy\,\psi)[t/x]=
\begin{cases}
 Qy\,\psi,&y=x,\\
 Qy\,(\psi[t/x]),&y\ne x\text{ and }y\notin\operatorname{Var}(t),\\
 Qz\,((\psi[z/y])[t/x]),&y\ne x\text{ and }y\in\operatorname{Var}(t),
\end{cases}
\]
where in the last clause \(z\) is the least variable not in
\[
 \operatorname{Var}(\psi)\cup\operatorname{Var}(t)\cup\set{x,y}.
\]
The recursion is on complexity, and it is well founded by a simultaneous
strong induction on complexity: the calls on \(\psi\) concern a formula of
lower complexity than \(Qy\,\psi\), and variable-for-variable
substitution preserves complexity---each clause preserves the outer
constructor---so in the last clause \(c(\psi[z/y])=c(\psi)<c(Qy\,\psi)\)
and the outer call is also at lower complexity.  Complexity preservation
for substitution by an arbitrary term is recorded as
\cref{lem:substitution-complexity} below.
\end{definition}

The last clause first renames the binder and then substitutes.  Its purpose is
not stylistic: it prevents a free variable of \(t\) from becoming bound.

\begin{lemma}[Substitution preserves formula complexity]
\label{lem:substitution-complexity}
For every formula \(\varphi\), term \(t\), and variable \(x\),
\[
  c(\varphi[t/x])=c(\varphi).
\]
\end{lemma}
\begin{proof}
Induct on \(c(\varphi)\).  Term substitution changes no logical constructor
in an atomic formula.  The connective clauses preserve the displayed outer
constructor and use the induction hypotheses.  At a quantifier, substitution
either stops, commutes with the quantifier, or renames its bound variable
before recurring on the matrix.  Bound-variable renaming changes no logical
constructor, so every quantifier case also preserves complexity.
\end{proof}

\begin{definition}[Alpha-equivalence]
The relation of \emph{alpha-equivalence}, written \(\equiv_\alpha\), is the
least equivalence relation on formulas that is a congruence for every formula
constructor and contains
\[
 Qx\,\varphi\equiv_\alpha Qy\,(\varphi[y/x])
\]
for \(Q\in\set{\forall,\exists}\) whenever
\(y\notin\operatorname{Var}(\varphi)\).  Thus a generating step renames one
bound variable and exactly the occurrences governed by its quantifier.
\end{definition}

Alpha-equivalence is syntactic and is not literal identity of strings.  It
will be shown below to preserve satisfaction.

\begin{example}[Avoiding capture]
Let \(\varphi=\forall y\,R(x,y)\) and substitute the term \(y\) for \(x\).
Direct textual replacement would produce \(\forall y\,R(y,y)\), incorrectly
binding the inserted occurrence.  If \(z\) is the least permitted fresh
variable, the defined result is
\[
  \varphi[y/x]=\forall z\,R(y,z).
\]
The free occurrence introduced by substitution remains free.
\end{example}

\begin{lemma}[Free variables under substitution]
\label{lem:substitution-fv}
For all terms \(r,t\) and every variable \(x\): if
\(x\notin\operatorname{Var}(r)\), then \(r[t/x]=r\); and
\[
 \operatorname{Var}(r[t/x])\subseteq
 (\operatorname{Var}(r)\setminus\set{x})\cup\operatorname{Var}(t).
\]
For every formula \(\varphi\), term \(t\), and variables \(x,y\):
\begin{enumerate}
  \item \(\FV(\varphi[t/x])\subseteq
        (\FV(\varphi)\setminus\set{x})\cup\operatorname{Var}(t)\); and
  \item if \(y\notin\operatorname{Var}(\varphi)\), then
        \(\FV(\varphi[y/x])=(\FV(\varphi)\setminus\set{x})\cup\set{y}\)
        when \(x\in\FV(\varphi)\), and
        \(\FV(\varphi[y/x])=\FV(\varphi)\) when
        \(x\notin\FV(\varphi)\).
\end{enumerate}
A term's variables are measured by \(\operatorname{Var}\), exactly as in
the \(\FV\) clauses for atomic formulas; no separate free-variable
operator on terms is defined in this document.
\textit{Status: complete proof.}
\end{lemma}
\begin{proof}
For the term claims, induct on \(r\).  A variable other than \(x\) is
unchanged, the variable \(x\) itself is replaced by \(t\), and the compound
clause applies the induction hypotheses argumentwise.

For the formula claims, proceed by strong induction on formula complexity
\(c\).  The atomic cases follow from the term claims and the \(\FV\)
clauses for atomic formulas; the \(\bottom\) case is immediate, and the
connective cases combine the induction hypotheses through the \(\FV\)
clauses.

For clause 2, the renaming clause of the substitution definition never
fires: every binder \(w\) of a subformula of \(\varphi\) lies in
\(\operatorname{Var}(\varphi)\), so \(w\ne y\), and the substitution is
textual replacement of the free occurrences of \(x\) by \(y\).  At a
quantifier \(Qw\,\psi\): if \(w=x\), substitution stops, \(x\) is not free,
and both sides equal \(\FV(Qw\,\psi)\); if \(w\ne x\), the substitution
commutes with the quantifier, and the induction hypothesis for \(\psi\)
passes through the clause
\(\FV(Qw\,(\psi[y/x]))=\FV(\psi[y/x])\setminus\set{w}\), using
\(w\ne y\) and the fact that \(x\in\FV(Qw\,\psi)\) iff
\(x\in\FV(\psi)\).

For clause 1 at a quantifier \(Qy\,\psi\): if \(y=x\), substitution stops
and \(x\notin\FV(Qy\,\psi)\); if \(y\ne x\) and
\(y\notin\operatorname{Var}(t)\), apply the induction hypothesis to
\(\psi\) and remove \(y\).  In the renaming case
\((Qy\,\psi)[t/x]=Qz\,((\psi[z/y])[t/x])\), apply the induction hypothesis
to \(\psi\) with \([z/y]\)---clause 2---and to \(\psi[z/y]\) with
\([t/x]\); the latter application is at the equal complexity
\(c(\psi[z/y])=c(\psi)<c(Qy\,\psi)\), licensed by
\cref{lem:substitution-complexity}.  Since
\(z\notin\operatorname{Var}(t)\cup\set{x}\), computing with the \(\FV\)
clauses gives
\[
 \FV\bigl(Qz\,((\psi[z/y])[t/x])\bigr)\subseteq
 (\FV(Qy\,\psi)\setminus\set{x})\cup\operatorname{Var}(t).
 \qedhere
\]
\end{proof}

\begin{lemma}
If \(x\notin\FV(\varphi)\), then
\(\varphi[t/x]\equiv_\alpha\varphi\).
\textit{Status: complete proof.}
\end{lemma}
\begin{proof}
Proceed by strong induction on formula complexity \(c\).  In the atomic
case, \(x\) occurs in no term of \(\varphi\), so every term is unchanged by
the term clause of \cref{lem:substitution-fv} and
\(\varphi[t/x]=\varphi\); the \(\bottom\) case is likewise immediate.  In
the connective cases, \(x\) is free in no immediate subformula; the
induction hypotheses and the congruence property of \(\equiv_\alpha\)
conclude.

At a quantifier \(Qy\,\psi\) with \(y=x\), substitution stops, and
reflexivity concludes.  If \(y\ne x\) and
\(y\notin\operatorname{Var}(t)\), then \(x\notin\FV(\psi)\); the induction
hypothesis gives \(\psi[t/x]\equiv_\alpha\psi\), and the quantifier
congruence concludes.  If \(y\ne x\) and \(y\in\operatorname{Var}(t)\),
the definition returns \(Qz\,((\psi[z/y])[t/x])\), where \(z\) is the
least variable outside
\(\operatorname{Var}(\psi)\cup\operatorname{Var}(t)\cup\set{x,y}\).  By
\cref{lem:substitution-fv}(2), \(x\notin\FV(\psi[z/y])\), because
\(x\notin\FV(\psi)\) and \(x\ne z\).  By
\cref{lem:substitution-complexity},
\(c(\psi[z/y])=c(\psi)<c(Qy\,\psi)\), so the strong induction hypothesis
applies to \(\psi[z/y]\) and gives
\((\psi[z/y])[t/x]\equiv_\alpha\psi[z/y]\).  The quantifier congruence
gives \(Qz\,((\psi[z/y])[t/x])\equiv_\alpha Qz\,(\psi[z/y])\); the
generating clause, licensed by \(z\notin\operatorname{Var}(\psi)\), gives
\(Qz\,(\psi[z/y])\equiv_\alpha Qy\,\psi\); transitivity concludes.
\end{proof}

